\section{Руководство по эксплуатации}
\label{sec:manual}

\subsection{Требования к окружению}

Перед запуском установите:

\begin{itemize}
    \item GHC версии \texttt{9.6} или новее;
    \item Stack версии \texttt{3} или новее;
    \item терминал с поддержкой UTF-8.
\end{itemize}

Для анализа подготовьте BDF-файл шрифта. Эталонный шрифт задаётся в конфигурации приложения; по умолчанию используется файл \texttt{data/Terminus 16v.bdf}.

\subsection{Сборка и запуск}

Склонируйте репозиторий:

\begin{minted}{text}
git clone https://github.com/si1og/SemenovIAFPCW.git
\end{minted}

Перейдите в каталог проекта:

\begin{minted}{text}
cd SemenovIAFPCW/font-analyzer
\end{minted}

Соберите программу:

\begin{minted}{text}
stack build
\end{minted}

Запустите программу:

\begin{minted}{text}
stack run
\end{minted}

Для проверки тестов используйте команду:

\begin{minted}{text}
stack test
\end{minted}

\subsection{Порядок работы}

После запуска выберите действие в меню. Доступные пункты приведены в таблице~\ref{tab:manual-menu}.

Постореть как выглядит меню в зпущенной программе можно постотреть на рисунке \ref{fig:img/main-menu.png}.


\begin{table}[H]
\centering
\caption{Пункты меню программы}
\label{tab:manual-menu}
\begin{tblr}{
    colspec={|Q[l,0.25\textwidth]|Q[l,0.65\textwidth]|},
    hlines,
    row{1}={font=\bfseries}
}
Пункт меню & Что сделать \\
\texttt{1} & Запустить анализ BDF-файла. \\
\texttt{2} & Включить или отключить отдельные метрики анализа. \\
\texttt{3} & Посмотреть текущие параметры анализа. \\
\texttt{0} & Выйти из программы. \\
\end{tblr}
\end{table}

Чтобы выполнить анализ, выберите пункт \texttt{1}, введите путь к BDF-файлу и затем путь для сохранения отчёта. Если оставить путь отчёта пустым, программа сохранит результат в файл \texttt{report.txt}.

\subsection{Настройка анализа}

Пункт меню \texttt{2} позволяет выбрать, какие метрики учитывать:

\begin{itemize}
    \item читаемость;
    \item пропорциональность;
    \item плотность;
    \item различимость.
\end{itemize}

На вопросы можно отвечать \texttt{y/n} или \texttt{да/нет}. Пустой ввод означает ответ по умолчанию.

Ввод метрик в прогркмме можно увидеть на рисунке \ref{fig:img/analysis-options.png}.

Пример диалога:

\begin{minted}{text}
анализировать читаемость? [y/n]
да
анализировать пропорциональность? [y/n]
нет
анализировать плотность? [y/n]
y
анализировать различимость? [y/n]

\end{minted}

В последнем вопросе пользователь оставил строку пустой, поэтому будет использовано значение по умолчанию.

\subsection{Формат входного файла}

На вход подаётся файл шрифта в формате BDF. Для корректной работы в описании глифа должны присутствовать строки:

\begin{itemize}
    \item \texttt{STARTCHAR} --- начало описания глифа;
    \item \texttt{ENCODING} --- код символа;
    \item \texttt{BITMAP} --- начало пиксельной матрицы;
    \item \texttt{ENDCHAR} --- конец описания глифа.
\end{itemize}

Строки пиксельной матрицы должны быть записаны в шестнадцатеричном виде.

\subsection{Результат работы}

После успешного анализа программа формирует текстовый отчёт. В отчёте содержатся:

\begin{itemize}
    \item общая информация о шрифте;
    \item описание выбранных метрик;
    \item таблица метрик по каждому глифу;
    \item список найденных аномалий;
    \item предложения замены из эталонного шрифта.
\end{itemize}

Пример начала сформированного отчёта приведён на рисунке~\ref{fig:img/report-summary.png}. Таблица метрик по глифам показана на рисунке~\ref{fig:img/report-glyph-metrics.png}. Пример списка найденных аномалий приведён на рисунке~\ref{fig:img/report-anomalies.png}, а предложения замены из эталонного шрифта показаны на рисунке~\ref{fig:img/report-replacements.png}.

События работы программы записываются в файл \texttt{coursework.log}.

\subsection{Обработка ошибок}

Если файл не найден, имеет неподдерживаемую структуру или отчёт не удалось записать, программа выводит сообщение об ошибке и возвращает управление пользователю. Информация об ошибке также записывается в журнал.

Примеры ошибок можно увидеть на рисунках \ref{fig:img/file-read-error.png} и \ref{fig:img/invalid-bdf-error.png}.

\newpage
